% Generated draft: requires mathematical and source review.
\documentclass[11pt]{article}
\usepackage[margin=25mm]{geometry}
\usepackage{fontspec}
\setmainfont{Latin Modern Roman}
\usepackage{amsmath,amssymb,mathtools,unicode-math}
\setmathfont{Latin Modern Math}
\usepackage{enumitem,needspace,xurl,hyperref}
\setlength{\parindent}{0pt}
\setlength{\parskip}{5pt}
\setlength{\emergencystretch}{3em}
\newcommand{\N}{\mathbb{N}}
\newcommand{\Z}{\mathbb{Z}}
\newcommand{\Q}{\mathbb{Q}}
\newcommand{\R}{\mathbb{R}}
\newcommand{\C}{\mathbb{C}}
\newcommand{\F}{\mathbb{F}}
\newcommand{\A}{\mathbb{A}}
\newcommand{\PP}{\mathbb{P}}
\newcommand{\E}{\mathbb{E}}
\newcommand{\eps}{\varepsilon}
\newcommand{\dd}{\,\mathrm{d}}
\begin{document}
\section*{UnsolvedMath ID: 701214}
\textbf{Source problem code:} catalogue entry
\subsection{Definitions and mathematical statement}
Let \(x\) be an indeterminate and let
\[
A=\Z[x,x^{-1}]
 =\left\{\sum_{k\in\Z}a_kx^k:\ a_k\in\Z,\ a_k=0\text{ for all but finitely many }k\right\}
\]
be the Laurent polynomial ring over the integers. Let \(GL_2(A)\) be the group of invertible \(2\times2\) matrices over \(A\), and define
\[
SL_2(A)=\{M\in GL_2(A):\det M=1\}.
\]
For \(t\in A\), the elementary matrices are
\[
e_{12}(t)=\begin{pmatrix}1&t\\0&1\end{pmatrix},\qquad
e_{21}(t)=\begin{pmatrix}1&0\\t&1\end{pmatrix}.
\]
Define
\[
E_2(A)=\bigl\langle e_{12}(t),e_{21}(t):t\in A\bigr\rangle\leq SL_2(A).
\]
The question attributed to L. Vaserstein is whether
\[
SL_2\bigl(\Z[x,x^{-1}]\bigr)=E_2\bigl(\Z[x,x^{-1}]\bigr).
\]
Equivalently, must every determinant-one \(2\times2\) matrix over \(A\) be a finite product of these elementary matrices? The supplied source describes this as a well-known problem and labels it open; that status is supplied metadata, not an independently verified literature assessment.
\subsection{Short English statement}
Is every determinant-one \(2\times2\) matrix over the integer Laurent polynomial ring \(\Z[x,x^{-1}]\) a finite product of elementary matrices? This is the supplied formulation of Kourovka Notebook Problem 16.22, attributed to L. Vaserstein.
\subsection{Sources}
\begin{itemize}
\item Local UnsolvedMath registry: permanent ID \(701214\), exact source problem code \texttt{catalogue entry}, title ``Kourovka Notebook Problem 16.22.'' The supplied record labels the problem open.
\item \emph{The Kourovka Notebook (Unsolved Problems in Group Theory)}, 16th issue (2006), Problem 16.22; proposed by L. Vaserstein. The supplied statement includes the qualification ``Well-known problem.''
\item The supplied background also identifies the 21st issue edition (2026), \texttt{arXiv:1401.0300v48}. The supplied collection description names E. I. Khukhro and V. D. Mazurov as editors. These bibliographic details and the status assessment have not been independently verified.
\end{itemize}
\end{document}
